\section{ANEXOS}
\addcontentsline{toc}{chapter}{ANEXOS}
\markboth{ANEXOS}{ANEXOS}

% --- ANEXO A: MATRIZ DE CONSISTENCIA ---
\section*{Anexo A: Matriz de Consistencia}

\begin{landscape}
\begin{table}[ht]
\centering
\caption{Matriz de Consistencia: Evaluación experimental de la precisión y rendimiento del código generado por inteligencia artificial en simulaciones físicas básicas.}
\label{tab:matriz_consistencia}
\scriptsize
\renewcommand{\arraystretch}{1.8}
\begin{tabularx}{\linewidth}{|X|X|X|X|X|}
\hline
\rowcolor[HTML]{EFEFEF} 
\centering \textbf{PROBLEMAS} & \centering \textbf{OBJETIVOS} & \centering \textbf{HIPÓTESIS} & \centering \textbf{VARIABLES Y DIMENSIONES} & \centering \textbf{METODOLOGÍA} \tabularnewline \hline

\textbf{Problema General:} \newline ¿Qué nivel de precisión física y rendimiento presenta el código generado por IA (GPT, Claude y Gemini) al aumentar la complejidad de las simulaciones? & 
\textbf{Objetivo General:} \newline Evaluar la precisión física y el rendimiento computacional del código generado por inteligencia artificial en escenarios de simulaciones mecánicas básicas. & 
\textbf{Hipótesis General:} \newline El código generado por IA presenta una degradación progresiva de la precisión física y un uso ineficiente de la CPU a medida que aumenta la complejidad del escenario físico. & 
\textbf{V. Independiente:} \newline - Herramienta de IA (GPT, Claude, Gemini). \newline - Escenario físico (Complejidad). \newline \newline \textbf{V. Dependiente 1:} \newline Error de Precisión. \newline \textit{(Dimensión: Desviación en unidades SI)}. \newline \newline \textbf{V. Dependiente 2:} \newline Rendimiento. \newline \textit{(Dimensión: Tiempo en ms y Memoria en MB)}. & 
\textbf{Tipo de Investigación:} \newline Aplicada y tecnológica. \newline \newline \textbf{Nivel:} \newline Descriptivo - Comparativo. \newline \newline \textbf{Diseño:} \newline Experimental comparativo transversal. \newline \newline \textbf{Método:} \newline Experimental. \newline \newline \textbf{Entorno:} \newline Arch Linux. \\ \hline

\textbf{Problemas Específicos:} \newline \newline 1. ¿Cuál es el error en caída libre (baja complejidad) respecto al modelo teórico analítico? \newline \newline 2. ¿Cómo afecta el aumento de la complejidad de las entidades al rendimiento y consumo de memoria? & 
\textbf{Objetivos Específicos:} \newline \newline 1. Determinar el error cuadrático medio en las variables de estado de caída libre generadas por la IA. \newline \newline 2. Medir el impacto del escalado de entidades y fricción en el tiempo de ejecución y el pico de consumo de memoria. & 
\textbf{Hipótesis Específicas:} \newline \newline 1. El error de precisión física será significativamente mayor en el código de la IA que en el modelo analítico de referencia. \newline \newline 2. A mayor complejidad del escenario físico, mayor será el consumo de recursos de CPU y memoria RAM del código generado. & 
& \\ \hline
\end{tabularx}
\end{table}
\end{landscape}

